% ECONOMICS_PROBLEM: {"id": "F13-419", "title": "Dynamic tariff incidence", "jel_code": "F13", "source_id": "OP-0342"}
% !TeX program = xelatex
\documentclass[11pt,a4paper]{article}
\usepackage[margin=23mm]{geometry}
\usepackage{fontspec}
\setmainfont{Latin Modern Roman}
\usepackage{amsmath,amssymb,mathtools}
\usepackage{xcolor,enumitem,booktabs,longtable,array,xurl,hyperref}
\definecolor{muted}{HTML}{53636C}
\hypersetup{colorlinks=true,urlcolor=blue,linkcolor=blue}
\setlength{\parindent}{0pt}
\setlength{\parskip}{5pt}
\newcommand{\R}{\mathbb R}
\newcommand{\E}{\mathbb E}
\newcommand{\N}{\mathbb N}
\newcommand{\ind}{\mathbf 1}
\newcommand{\Var}{\operatorname{Var}}
\newcommand{\Cov}{\operatorname{Cov}}
\newcommand{\supp}{\operatorname{supp}}
\DeclareMathOperator*{\argmin}{arg\,min}
\DeclareMathOperator*{\argmax}{arg\,max}
\newcommand{\entrymeta}[3]{{\sffamily\small\color{muted}\textbf{\texttt{#1}}\quad|\quad #2\par #3\par}}
\newcommand{\Problemref}[1]{\texttt{#1}}
\newcommand{\JELref}[2]{\texttt{#2} (JEL \texttt{#1})}
\begin{document}
\section*{Dynamic tariff incidence}
\textbf{Problem ID:} \texttt{F13-419}\quad \textbf{JEL:} \texttt{F13}\par
% BEGIN ECONOMICS STATEMENT
\label{problem:OP-0342}
{\sffamily\small\color{muted}Original subject group: Trade, globalization and geoeconomics\par}
\label{definitions:problem:30007624}
\textbf{Audit classification:} Published research agenda. The 2026 primary draft explicitly lists future work on dynamic and distributional tariff effects. The stage-specific potential-price formula is editorial; already estimated short-run results are distinguished.
\subsection*{Definitions and precise research target}
Fix a product-level tariff intervention increasing the statutory ad valorem rate from \(\tau_0\) to \(\tau_1\), with \(0\leq\tau_0<\tau_1\), at date \(t\). Let \(P_{k,t}\) denote the price at supply-chain stage \(k\) and date \(t\). Track foreign export prices, tariff-inclusive import prices, domestic producer prices, retail prices, quantities, wages and profits. For each stage \(k\) and horizon \(h\), define pass-through
\[\rho_k(h)=\frac{E[\log P_{k,t+h}(do(\tau_1))-\log P_{k,t+h}(do(\tau_0))]}{\log(1+\tau_1)-\log(1+\tau_0)}.\]
The task is to identify these dynamic responses and allocate the welfare burden across foreign producers, importers, retailers, workers and consumers, counting tariff revenue consistently. Specify product composition, exchange-rate treatment and the counterfactual policy path, including retaliation. Endogenous sourcing changes mean price indexes must distinguish within-product repricing from entry, exit and quality substitution. Anticipatory stockpiling and inventory drawdown must be modeled or measured before interpreting delayed retail effects. Establish conditions for a valid comparison group when tariff assignments target politically or economically distinct products. Report identified sets where export margins or profits are unobserved. High import-price pass-through and low contemporaneous consumer-price pass-through are compatible and do not resolve how burdens evolve after firms exhaust inventories or reorganize supply chains.

\textbf{Additional formal definitions and scope (editorial).} Fix the same product varieties and quality units at each supply-chain stage, or state a price-index formula accounting for variety entry and exit. \(P_{k,t+h}(a)>0\) is the potential stage-\(k\) transaction price; horizons are \(h\in\mathbb N_0\). The denominator is a change in the logarithmic tariff-inclusive factor, hence \(\rho_k\) is a finite-change elasticity relative to that factor. Foreign export and tariff-inclusive import prices require common currency conversion and the same delivery terms. Define inventory identities, contract lags and tariff collection dates. Welfare burden includes producer/consumer changes and the domestic government budget; tariff revenue is a transfer from payers, while compliance and distortion are resource costs. Indirect consumer taxes and retaliation must either change under the specified intervention or remain fixed by assumption.

For the identification part, declare an admissible class \(\Theta\) of structural models, including preferences, technologies, shock laws, measurement/sampling rules and any equilibrium selector. If \(O\) denotes the complete observable history and \(T(\theta)\) the target defined above, the identified set is \(\mathcal I_T=\{T(\theta):\theta\in\Theta,\ \mathcal L_\theta(O)=\mathcal L_{\rm obs}(O)\}\); point identification means this set is a singleton. A \(do(a)\) comparison replaces stated policy/technology equations while fixing the declared exogenous law and re-solving the remaining equations. These definitions do not assert that an identification theorem holds without further restrictions.
\subsection*{Variants and assumption changes}
\begin{enumerate}
\item \textbf{Short-run fixed sourcing (source-motivated).} Hold supplier contracts and product coverage fixed, measure tariff-inclusive import prices and inventory-mediated retail pass-through.
\item \textbf{Dynamic sourcing and incidence (source-motivated extension).} Let inventories, entry, contracts and retaliation adjust; identify burden across agents and terms-of-trade effects over a specified policy path.
\end{enumerate}
\subsection*{References and source audit}
\begin{enumerate}
\item Brookings cover verifies Fajgelbaum and Khandelwal, March 2026 conference draft; NBER indexed metadata independently verifies WP 35064 and DOI. Locator: Section 7, printed p.49 (PDF p.50), explicit further-research list. Support: Dynamic and distributional tariff agenda. Full conference draft accessible. URL: \url{https://www.brookings.edu/wp-content/uploads/2026/03/1_Fajgelbaum-Khandelwal_unembargoed.pdf}.
\end{enumerate}
\begin{enumerate}\raggedright
\item\hypertarget{definitions:source:30007624:1}{} Pablo D. Fajgelbaum; Amit Khandelwal. 2026. \emph{Tariffs in 2025: Short-Run Impacts on the U.S. Economy}. Section 7, printed p.49 (PDF p.50), explicit further-research list.
\url{https://www.brookings.edu/wp-content/uploads/2026/03/1_Fajgelbaum-Khandelwal_unembargoed.pdf}.
DOI: \url{https://doi.org/10.3386/w35064}.
\end{enumerate}

\subsection*{Common conventions: Source mathematical conventions}

These conventions apply to each entry unless explicitly replaced.

\textbf{Models and identification.} A primitive model \(m\) specifies all probability laws, feasible choices, information, equilibrium and measurement rules used by its target. The admissible class \(\mathcal M\) is fixed independently of outcomes. If \(P_O(m)\) is its observable probability law and \(\tau(m)\) the target, its identified set at observable law \(P\) is
\[
 \mathcal I_\tau(P)=\{\tau(m):m\in\mathcal M,\ P_O(m)=P\}.
\]
Point identification means that this set is a singleton. An empty set signals incompatibility of data and assumptions. Coordinate bounds need not describe the joint set. Bounds are sharp only if every retained target is attainable or arbitrarily approachable by compatible models. Finite bounds require explicit restrictions on unobserved quantities; unrestricted confounding, hidden wealth, extrapolation or arbitrary equilibrium selection can yield uninformative sets.

\textbf{Causal interventions.} A potential outcome \(Y_i(p)\) is the outcome under a complete policy regime \(p\), including timing, financing, information and market spillovers. The target population and units are fixed before treatment. With interference, use the complete assignment vector or a justified exposure mapping; individual no-interference notation alone cannot identify an economy-wide intervention. A difference of mean potential outcomes does not require the cross-world joint distribution. Conditioning on post-treatment migration, survival, enrollment or employment changes the estimand and needs separate justification.

\textbf{Statistical statements.} An estimator, confidence set, forecast or test specifies a sampling law, model class and loss/coverage criterion. Uniform \((1-\alpha)\)-coverage means \(\inf_{m\in\mathcal M}\Pr_m\{\tau(m)\in C_n\}\geq1-\alpha\), or an explicitly asymptotic version. The whole parameter space trivially has coverage, so informativeness requires a stated width, risk or power objective. A forecasting improvement is relative to a named benchmark, information set and evaluation population.

\textbf{Equilibrium and optimization.} An equilibrium comprises feasible plans, optimality, beliefs where needed, budget constraints and all market/resource clearing. The primitives or a stated benchmark define each correspondence \(\mathcal E(p;m)\). Nonemptiness is assumed or proved, never inferred just from compact policy sets. A selected-equilibrium target uses a declared selection rule; a robust target quantifies over all permitted equilibria. Use suprema or infima unless attainment is established. An \(\varepsilon\)-optimal policy is feasible and within \(\varepsilon\) of the finite optimum value. Report an empty feasible set separately.

\textbf{Welfare and accounting.} Normative utility, interpersonal normalization, social weights, numeraire, discounting and the use of public receipts are primitives. Behavioral utility need not coincide with the welfare criterion. Household welfare includes ownership income and rents once; adding profits again to compensating variation generally double-counts them. A monetary equivalent is defined by an explicit transfer and price convention, with monotonicity and range conditions. Average effects, mechanism contributions, transfers and welfare losses are different targets. Infinite sums require convergence and dynamic feasible plans require terminal or no-Ponzi conditions.

\textbf{Source versions.} A working-paper date, revision date and journal publication year are distinguished. A DOI for a journal version is not automatically the DOI of an earlier manuscript. Locators refer to the stated version and printed pages where specified. A metadata-only check does not verify a claim about a full-text conclusion. Every entry's source subsection narrows the provenance claim accordingly.
% END ECONOMICS STATEMENT
\end{document}
